% O Overleaf já inclui IEEEtran; use este arquivo pelo main.tex na raiz do projeto.
\documentclass[conference]{IEEEtran}
\usepackage[T1]{fontenc}
\usepackage[utf8]{inputenc}
\usepackage[brazilian]{babel}
\usepackage{graphicx}
\usepackage{booktabs}
\usepackage{url}
\usepackage{float}

\title{Autoescalamento Horizontal de uma Aplicação Web em Kubernetes: comparação entre Minikube e Amazon EKS}

\author{\IEEEauthorblockN{Mario Marszalek}
\IEEEauthorblockA{PSI5120 -- Tópicos em Computação em Nuvem\\
Escola Politécnica -- Universidade de São Paulo\\
NUSP: 11910490 -- mariomarszalek@usp.br}}

\begin{document}
\maketitle

\begin{abstract}
Este trabalho implementa e avalia o autoescalamento horizontal de uma aplicação web em Kubernetes. Foram feitas duas implantações do mesmo sistema: uma local, usando Minikube, e outra gerenciada, usando Amazon EKS. A aplicação em PHP recebe uma carga de CPU controlada por parâmetro; o Horizontal Pod Autoscaler (HPA) usa CPU média, com alvo de 50\%, para variar entre uma e cinco réplicas. No EKS, a utilização registrada passou de 1\% para 498\% do \textit{request} de CPU, enquanto a meta do HPA era 50\%; a captura do pico mostra quatro réplicas ativas. O experimento também mostra a diferença prática entre o ambiente local, rápido para iterar, e o ambiente gerenciado, que exige mais configuração, credenciais e cuidado com custos.
\end{abstract}

\begin{IEEEkeywords}
Kubernetes, Horizontal Pod Autoscaler, HPA, Minikube, Amazon EKS, computação em nuvem.
\end{IEEEkeywords}

\section{Introdução}
Aplicações web têm demanda variável. Manter um número fixo de processos pode deixar o sistema lento em picos ou gastar recurso à toa quando não há tráfego. Em Kubernetes, o Horizontal Pod Autoscaler ajusta a quantidade de réplicas de um recurso escalável, como um \textit{Deployment}, a partir de métricas~\cite{k8s-hpa}. Neste trabalho foi usado o HPA baseado em utilização média de CPU.

O objetivo não foi montar uma aplicação complexa. A ideia foi montar um caso pequeno que desse para medir: uma aplicação web, um \textit{Service}, um \textit{Deployment}, um gerador de carga e o HPA. O mesmo desenho foi executado localmente no Minikube e na nuvem no Amazon EKS. Assim, foi possível observar tanto o comportamento do escalonamento quanto o custo operacional de subir a infraestrutura.

\section{Arquitetura da solução}
A Figura~\ref{fig:arquitetura} resume o fluxo. O \textit{Job} de carga faz requisições para o \textit{Service}. O \textit{Service} distribui as requisições entre os \textit{pods} da aplicação. O Metrics Server fornece as métricas de CPU e o HPA altera o número de réplicas do \textit{Deployment}.

\begin{figure}[H]
\centering
\fbox{\parbox{0.43\textwidth}{\centering
Gerador de carga $\rightarrow$ Service \texttt{web-hpa} $\rightarrow$ Deployment \texttt{web-hpa} (1--5 pods)\\[2mm]
Metrics Server $\rightarrow$ HPA (CPU média, meta 50\%)}}
\caption{Topologia lógica usada nas duas implantações.}
\label{fig:arquitetura}
\end{figure}

A aplicação foi empacotada em uma imagem PHP/Apache. Ela expõe uma rota HTTP e aceita o parâmetro \texttt{work}, que executa um cálculo determinístico para elevar o uso de CPU. Cada contêiner solicita 100m de CPU e 64Mi de memória, com limites de 500m e 128Mi. Esses \textit{requests} são importantes: a métrica de utilização do HPA é calculada em relação ao \textit{request}, não ao limite.

O HPA foi configurado com mínimo de uma e máximo de cinco réplicas. Para subida, foi usada uma janela de estabilização de zero segundo e uma política de até 100\% a cada 15 segundos. Para descida, foi usada janela de 60 segundos. Isso evita oscilações logo depois que a carga para.

\section{Implantação local com Minikube}
O ambiente local foi criado com Minikube usando o driver Docker. O Metrics Server foi habilitado como \textit{addon}, a imagem foi construída localmente e carregada para o ambiente Minikube. Em seguida foram aplicados \textit{namespace}, \textit{Deployment}, \textit{Service} e HPA.

Antes da carga, havia uma réplica em execução e a utilização observada era 1\% frente à meta de 50\%, como mostra a Figura~\ref{fig:mini-antes}. Durante a carga, o HPA detectou utilização acima da meta e aumentou as réplicas (Figura~\ref{fig:mini-durante}). Após remover o \textit{Job}, a utilização voltou a um valor baixo (Figura~\ref{fig:mini-pos}). A permanência temporária de réplicas extras no pós-teste é esperada por causa da janela de estabilização.

\begin{figure}[H]
\centering
\includegraphics[width=\columnwidth]{minikube-hpa/evidence/04-ciclo-continuo-antes.png}
\caption{Minikube antes da carga: uma réplica e CPU em 1\% da referência do HPA.}
\label{fig:mini-antes}
\end{figure}

\begin{figure}[H]
\centering
\includegraphics[width=\columnwidth]{minikube-hpa/evidence/05-ciclo-continuo-durante.png}
\caption{Minikube durante a carga, com HPA e pods adicionais.}
\label{fig:mini-durante}
\end{figure}

\begin{figure}[H]
\centering
\includegraphics[width=\columnwidth]{minikube-hpa/evidence/06-ciclo-continuo-pos.png}
\caption{Minikube após remover a carga: CPU retorna a 1\%; a redução de réplicas respeita a janela de estabilização.}
\label{fig:mini-pos}
\end{figure}

\section{Implantação em Amazon EKS}
Na nuvem, foi criado o cluster \texttt{psi5120-hpa} na região \texttt{sa-east-1} (São Paulo). A imagem da aplicação foi publicada no Amazon ECR e o \textit{Deployment} do EKS referencia essa imagem. Foi utilizado um único \textit{managed node group}, sem NAT Gateway, para reduzir o custo do experimento. O Metrics Server foi instalado como \textit{add-on} do EKS~\cite{eks-hpa}.

Houve uma restrição relevante da conta: o tipo \texttt{t3.medium} foi rejeitado por não ser elegível ao Free Tier. O teste foi então executado com \texttt{t3.small}, que a própria conta listou como elegível. Esse desvio é relevante porque mostra que a escolha de máquina não pode ser assumida apenas pelo tutorial; ela depende das regras e créditos disponíveis na conta.

A Figura~\ref{fig:eks-console} registra o cluster ativo no Console da AWS, com Kubernetes 1.34 e região São Paulo. A infraestrutura foi mantida pequena e o \textit{namespace} de teste foi removido logo após a coleta das evidências. Os \textit{node groups} foram colocados em exclusão para evitar custo residual, como mostra a Figura~\ref{fig:eks-cleanup}.

\begin{figure}[H]
\centering
\includegraphics[width=\columnwidth]{minikube-hpa/evidence-eks/04-console-eks-cluster-active.png}
\caption{Console da AWS: cluster EKS ativo antes da limpeza.}
\label{fig:eks-console}
\end{figure}

\begin{figure}[H]
\centering
\includegraphics[width=\columnwidth]{minikube-hpa/evidence-eks/06-console-eks-cleanup.png}
\caption{Console da AWS durante a limpeza: grupos de nós em exclusão após os testes.}
\label{fig:eks-cleanup}
\end{figure}

\section{Metodologia de teste}
O teste seguiu três momentos: (i) confirmar o estado de repouso com \texttt{kubectl top pods} e \texttt{kubectl get hpa}; (ii) iniciar o \textit{Job} de carga; e (iii) remover o \textit{Job} e observar o retorno das métricas. As capturas foram feitas com os comandos \texttt{kubectl get hpa,pods} e \texttt{kubectl top pods}. Não foram usados valores estimados nas evidências: os números apresentados são os observados no cluster.

O Job \texttt{web-hpa-load} usa a imagem \texttt{busybox:1.36}. Seu comando inicia 20 processos concorrentes. Cada processo executa requisições HTTP com \texttt{wget} para o Service interno \texttt{web-hpa.psi5120-hpa.svc.cluster.local}, com o parâmetro \texttt{work=75000}, durante até 300 segundos. A carga fica dentro do cluster, portanto não depende de Load Balancer público nem de tráfego externo. Para encerrar o ensaio antes dos cinco minutos, o Job é removido explicitamente.

O parâmetro \texttt{work} é consumido pela aplicação PHP e provoca cálculo de CPU antes da resposta. O Deployment começa com uma réplica. Cada pod solicita 100m de CPU e 64Mi de memória, e tem limites de 500m e 128Mi. Como a meta do HPA é 50\% de CPU, uma média próxima de 50m por pod já é suficiente para iniciar uma decisão de subida. A \textit{readiness probe} HTTP usa atraso inicial de três segundos, período de cinco segundos e até seis falhas; assim, um novo pod só passa a receber tráfego depois de responder adequadamente.

O mesmo gerador de carga e a mesma política de HPA foram usados nas duas implantações. A diferença entre elas ficou restrita ao ambiente de execução e à origem da imagem: no Minikube a imagem foi carregada localmente, enquanto no EKS ela foi obtida do ECR. Isso reduz uma variável de confusão na comparação, pois a regra de escalonamento e o perfil de carga permaneceram os mesmos.

\section{Disponibilidade do repositório}
O código-fonte, os manifestos, os roteiros de execução e as evidências usados neste trabalho estão disponíveis em \url{https://github.com/mmarsz/psi5120-hpa/}. O repositório contém o Dockerfile e a aplicação PHP, os manifestos de \textit{namespace}, Deployment, Service, HPA e Job de carga, além da configuração \texttt{eksctl} usada para o EKS.

As evidências estão separadas por ambiente. A pasta \texttt{minikube-hpa/evidence} contém as capturas do ciclo contínuo local, incluindo os estados antes, durante e após a carga. A pasta \texttt{minikube-hpa/evidence-eks} contém as capturas e transcritos da implantação gerenciada, incluindo o Console da AWS. Essa organização permite relacionar cada resultado discutido no artigo a um arquivo verificável.

O repositório também registra decisões necessárias para reproduzir o laboratório: imagem local no Minikube, imagem no ECR para EKS, alvo de CPU de 50\%, limites de uma a cinco réplicas e ausência de NAT Gateway. O acesso pode ser liberado aos docentes para verificação da entrega, conforme solicitado no enunciado.

No EKS, antes da carga, a única réplica consumia 1m de CPU, isto é, 1\% do \textit{request} de referência (Figura~\ref{fig:eks-antes}). Durante a carga, o HPA apresentou 498\%/50\% e quatro réplicas ativas (Figura~\ref{fig:eks-durante}). Após remover o Job, a métrica retornou a 1\%, enquanto a redução de réplicas respeitava a janela configurada (Figura~\ref{fig:eks-pos}). A carga também aparecia como um pod separado, consumindo 674m de CPU na observação do pico.

\begin{figure}[H]
\centering
\includegraphics[width=\columnwidth]{minikube-hpa/evidence-eks/01-antes-carga.png}
\caption{EKS antes da carga: HPA em 1\%/50\% e uma réplica.}
\label{fig:eks-antes}
\end{figure}

\begin{figure}[H]
\centering
\includegraphics[width=\columnwidth]{minikube-hpa/evidence-eks/02-durante-carga.png}
\caption{EKS durante a carga: HPA em 498\%/50\% e escalonamento em andamento.}
\label{fig:eks-durante}
\end{figure}

\begin{figure}[H]
\centering
\includegraphics[width=\columnwidth]{minikube-hpa/evidence-eks/03-pos-carga.png}
\caption{EKS após remover a carga: CPU retorna a 1\%; a descida respeita a janela do HPA.}
\label{fig:eks-pos}
\end{figure}

\section{Resultados e discussão}
A Tabela~\ref{tab:comparacao} reúne os resultados práticos. A maior diferença está na operação. No Minikube, o ciclo de teste foi direto: criar o cluster, carregar a imagem e aplicar manifestos. No EKS, foi necessário configurar a imagem no ECR, criar o control plane, criar um grupo de nós, atualizar o \textit{kubeconfig} e instalar o Metrics Server. Em troca, o EKS entrega um ambiente gerenciado e mais próximo do uso real em nuvem.

\begin{table}[H]
\caption{Comparação observada entre os ambientes}
\label{tab:comparacao}
\centering
\begin{tabular}{p{0.16\columnwidth}p{0.30\columnwidth}p{0.30\columnwidth}}
\toprule
Item & Minikube & Amazon EKS \\
\midrule
Execução & Local, com Docker e imagem carregada no cluster & Gerenciada, com ECR, EKS e nó EC2 \\
Métrica & Metrics Server como addon & Metrics Server como add-on EKS \\
Escala configurada & 1 a 5 pods & 1 a 5 pods \\
Evidência de pico & HPA e pods aumentados & 498\%/50\% e 4 réplicas ativas no registro \\
Gestão & Menos etapas, ideal para desenvolvimento & Mais etapas, IAM e controle de custo \\
\bottomrule
\end{tabular}
\end{table}

Sobre custo, não foi calculado um valor em reais neste artigo porque o laboratório usou créditos estudantis e teve execução curta. Ainda assim, o desenho tomou decisões para reduzir gasto: um nó pequeno, sem NAT Gateway, sem Load Balancer e remoção da aplicação após o teste. A criação de EKS não deve ser deixada ativa depois da entrega, pois o control plane e os recursos associados podem gerar cobrança.

\section{Tempo de reação e custo operacional}
No Minikube foi registrado um ciclo contínuo com horário. Às 23:25:07, antes da carga, o HPA indicava 1\%/50\% e uma réplica. Entre essa observação e a leitura de 23:26:04, o Job foi iniciado; 57 segundos depois, a utilização estava em 67\%/50\% e o HPA já indicava duas réplicas. Esses registros não medem latência HTTP; eles medem o intervalo operacional entre iniciar a carga, coletar métricas e observar o escalonamento.

Após remover o Job, a leitura de 23:28:47 voltou a 1\%/50\%. As cinco réplicas ainda apareciam no instante dessa captura porque a política de descida usa janela de estabilização de 60 segundos. Portanto, a redução da métrica ocorreu antes da redução da quantidade de pods, comportamento compatível com a configuração do HPA.

No EKS, a captura de repouso mostra 1\%/50\% e uma réplica; durante a carga, mostra 498\%/50\% e quatro réplicas ativas. A infraestrutura em nuvem exigiu mais etapas que Minikube: imagem no ECR, criação do control plane, grupo de nós, atualização do \textit{kubeconfig} e instalação do Metrics Server. Para limitar custo, foi usado um único nó \texttt{t3.small}, sem NAT Gateway e sem Load Balancer. Após os testes, o namespace foi removido, os grupos de nós foram colocados em exclusão e a exclusão do cluster foi iniciada.

\section{Conclusão}
O objetivo foi atingido nas duas plataformas: a aplicação web foi implantada, o Metrics Server forneceu métricas e o HPA reagiu à carga de CPU. No EKS, a observação de 498\% frente à meta de 50\% e o aumento de uma para múltiplas réplicas deixam o comportamento visível. O retorno para 1\% após remover a carga confirma que o gerador de carga era a causa do aumento.

Minikube foi mais simples para validar a lógica e ajustar manifestos. EKS exigiu mais trabalho de infraestrutura, especialmente credenciais, imagem no ECR e compatibilidade de tipo de instância, mas forneceu a validação em ambiente de nuvem gerenciada solicitada pelo enunciado. Como trabalho futuro, seria interessante repetir o ensaio com mais nós, medir latência de resposta e registrar custos detalhados por serviço.

\appendices
\section{Roteiro de reprodução}
Este apêndice deixa o caminho usado no laboratório explícito. No Minikube, o cluster foi iniciado com driver Docker, quatro CPUs e 6GiB de memória. O Metrics Server foi habilitado, a imagem foi construída localmente e carregada no cluster. Depois foram aplicados \texttt{namespace.yaml}, \texttt{deployment.yaml}, \texttt{service.yaml} e \texttt{hpa.yaml}. A verificação inicial foi feita com \texttt{kubectl get deploy,pods,svc,hpa} e \texttt{kubectl top pods}.

No EKS, o arquivo \texttt{eksctl-cluster.yaml} definiu a região São Paulo, um nó gerenciado, acesso público ao endpoint e ausência de NAT Gateway. Após o cluster ficar ativo, o \textit{kubeconfig} foi atualizado com \texttt{aws eks update-kubeconfig}. A imagem já publicada no ECR foi referenciada por \texttt{deployment-eks.yaml}. O Metrics Server foi instalado como add-on e a aplicação foi aplicada com os mesmos manifestos lógicos usados localmente.

O gerador de carga é um \textit{Job} separado. Isso foi útil porque iniciar e encerrar a carga não exige alterar o \textit{Deployment} principal. Durante o ensaio, os comandos de observação ficaram separados da aplicação: \texttt{kubectl get hpa}, \texttt{kubectl get pods} e \texttt{kubectl top pods}. No final, o Job foi removido com \texttt{kubectl delete job web-hpa-load}.

\section{Leitura das evidências}
As capturas não devem ser lidas como medição de desempenho absoluto da AWS. O objetivo do ensaio foi verificar a reação do HPA. O valor de 498\%/50\% significa que a utilização média calculada pelo HPA estava muito acima da meta configurada; por isso, o controlador aumentou as réplicas. As quatro réplicas ativas observadas na captura mostram que a ação não ocorre como uma troca atômica: pods precisam ser criados, agendados e passar pela \textit{readiness probe}.

No pós-carga, a CPU retornou para aproximadamente 1m por pod. Ainda assim, o número de réplicas não caiu imediatamente. Isso é comportamento esperado e foi configurado de propósito: a janela de estabilização de descida de 60 segundos impede que pequenas oscilações de tráfego façam o sistema escalar para cima e para baixo o tempo todo.

\section{Limitações e decisões}
O teste usa uma única máquina no EKS. Logo, ele prova escalonamento de pods, não escalonamento de nós. Quando o HPA pede mais pods do que cabem no nó, novos pods podem esperar por capacidade. Em um ambiente de produção, seria possível combinar HPA com Cluster Autoscaler ou Karpenter e definir mais de um nó. Isso não foi feito porque o foco do trabalho era HPA e porque o laboratório tinha orçamento limitado.

Também não foi criado Load Balancer público. O Service é do tipo ClusterIP e o tráfego de teste vem de dentro do cluster. Essa escolha reduz custo e deixa o experimento reproduzível sem expor a aplicação à Internet. Por fim, o valor financeiro não foi estimado como se fosse uma medição precisa; o artigo registra apenas as escolhas de contenção de custo e a limpeza iniciada depois das evidências.

\section{Cronologia do experimento em EKS}
A execução em nuvem foi feita em etapas para evitar confundir falha de infraestrutura com falha do HPA. Primeiro, o control plane foi criado e sua situação foi conferida como \textit{Active}. Segundo, o nó gerenciado foi criado e confirmado como \textit{Ready} no Kubernetes. Terceiro, a imagem no ECR, o Metrics Server e os manifestos foram validados. Somente então o Job de carga foi iniciado.

Durante a criação do nó ocorreu uma rejeição de tipo de instância pelo Free Tier. A tentativa foi interrompida, o grupo em erro foi colocado em exclusão e foi usado um novo grupo com \texttt{t3.small}. Essa sequência foi preservada nas evidências do Console de limpeza. O ponto importante é que o teste final não dependeu de um recurso em estado de erro: a aplicação foi executada em um nó que apareceu como \textit{Ready} e o cluster foi observado ativo pelo Console e pelo \texttt{kubectl}.

O HPA ficou inicialmente com uma réplica. O primeiro ciclo de métricas reportou 1\%/50\% e, no pico registrado, 498\%/50\%. Esse intervalo ilustra a natureza periódica do controlador: a escala é decidida depois de coletar métricas e os novos pods levam algum tempo para ficarem disponíveis. Por isso, o artigo usa as observações em sequência e não um único número isolado.

\section{Rastreabilidade dos arquivos}
O repositório foi organizado para que a correção não dependa de comandos escondidos. O \texttt{Dockerfile} e \texttt{app/index.php} definem a aplicação. A pasta \texttt{manifests} contém os recursos Kubernetes; há uma versão local do Deployment e outra para EKS, que troca apenas a origem da imagem pelo ECR. O arquivo \texttt{eksctl-cluster.yaml} registra a topologia do cluster gerenciado e \texttt{cleanup-eks.sh} concentra a remoção da infraestrutura.

As imagens e os respectivos transcritos de terminal ficam em duas pastas. \texttt{evidence} contém o Minikube, e \texttt{evidence-eks} contém o EKS. Os arquivos são nomeados pelo instante do teste; as referências no texto usam a numeração gerada pelo LaTeX. A captura do Console do cluster ativo está na Figura~\ref{fig:eks-console}. Essa separação permite que cada afirmação do texto seja conferida em um artefato do repositório.

Para reproduzir o experimento, deve-se manter os comentários dos manifestos e ajustar apenas credenciais, nome de imagem no ECR e dados da conta. Nenhuma credencial é armazenada no repositório. Esse cuidado é especialmente importante porque o trabalho pede um link que pode ser compartilhado com docentes.

\section{Ameaças à validade do experimento}
Os resultados devem ser interpretados como validação funcional do HPA, e não como \textit{benchmark} de desempenho entre plataformas. Minikube e EKS foram executados em recursos diferentes, com camadas de virtualização, rede e disponibilidade de CPU distintas. Por esse motivo, percentuais de CPU e quantidade de réplicas comprovam a reação do controlador em cada ambiente, mas não permitem concluir que uma plataforma é mais rápida que a outra.

O experimento no EKS usou apenas um nó \texttt{t3.small}. Assim, o HPA pode solicitar novos pods sem que isso implique criação automática de novos nós. Essa decisão foi deliberada para manter o escopo no autoescalamento de pods e limitar o uso de créditos estudantis. Em uma implantação de produção, a ausência de capacidade no nó precisaria ser tratada por um mecanismo complementar de escalonamento de nós.

Por fim, as métricas são amostras do Metrics Server e do ciclo de reconciliação do HPA. A captura durante a carga registra um estado observado, não uma série temporal completa. Mesmo com essa limitação, as evidências mostram a sequência essencial: repouso com uma réplica, CPU acima da meta com novas réplicas e retorno da CPU a 1\% depois da remoção do Job.

\clearpage
\begin{thebibliography}{9}
\bibitem{k8s-hpa} Kubernetes, ``Horizontal Pod Autoscaling Walkthrough.'' Disponível em: \url{https://kubernetes.io/docs/tasks/run-application/horizontal-pod-autoscale-walkthrough/}. Acesso em: 31 ago. 2026.
\bibitem{eks-hpa} Amazon Web Services, ``Scale pod deployments with Horizontal Pod Autoscaler.'' Disponível em: \url{https://docs.aws.amazon.com/eks/latest/userguide/horizontal-pod-autoscaler.html}. Acesso em: 31 ago. 2026.
\end{thebibliography}
\end{document}
